% SPDX-License-Identifier: Apache-2.0
\section{The Machine That Restarts Itself}
\label{sec:x-wdog}

A program that hangs stays hung until somebody presses the button, and
on a board nobody is standing next to, nobody does. So the hardware
counts, and software has to keep saying that it is still there.

\code{Wdog} is that counter in six words. A timeout, the count left of
it, a control word, a status, the window's sill, and the register a
refresh is written to. When the count reaches zero it asks the system
control block of Section~\ref{sec:x-syscon} for a reset, which records
the watchdog as the cause, so a program that restarts can read why it
did.

\subsection{Why A Refresh Carries A Key}

A program that has lost its way is not a program that has stopped. It
is a program writing whatever it finds over whatever it reaches, and
it will eventually reach the watchdog. A refresh that any write
performs is a refresh such a program keeps doing by accident, and a
watchdog it cannot fail is no watchdog.

So a refresh carries a word, and a write of anything else is not
ignored: it is a failure, and it asks for the reset at once. The
example writes a zero and watches the status.

\subsection{The Window, And What It Catches}

The plain countdown catches a program that stopped. It does not catch
a program looping tightly on its own refresh, which is just as stuck
and looks from here like the healthiest program on the machine.

In window mode a refresh is too early while the count is still above
the sill, and too early is a failure like a wrong key. A program then
has to refresh in the second half of its timeout, which it can only do
if it is keeping time rather than spinning. The example sets the sill
at half the timeout and refreshes twice in quick succession, which is
what a spinning program looks like.

\subsection{Warning First}

A timeout can mean the machine is broken, and it can mean that
something took longer than it usually does. With warning on, the first
timeout raises an interrupt and reloads the count instead of resetting,
which gives a program that is merely slow one more timeout to say so.
The second one resets. The run below takes both. The ticks it prints
are when the program read the status rather than when the watchdog
moved, since a read over this bus costs about twenty six cycles; the
figure is the moment itself.

\subsection{The Lock}

A watchdog software can turn off is a watchdog that gets turned off,
usually while somebody is debugging something else. The lock bit sets
and does not clear, and after it is set the control, the timeout and
the sill are read only until the next reset. The example locks the
watchdog on and then writes a zero over the control word, which
changes nothing.

\begin{figure}[htbp]
\centering
\input{wdog_timing.tex}
\caption{The second timeout: the count falling to zero with the
warning already standing on \code{irq}, the reset asked for on
\code{rst\_req} and held while \code{hold} counts down, the status going
from warned to warned and failed, and the count reloaded so that the
request ends rather than standing for ever.}
\label{fig:x-wdog}
\end{figure}

\lstinputlisting[caption={\code{ex\_wdog.rs}. Run with \code{bazel run //lib/examples:ex\_wdog}.},label={lst:x-wdog}]{lib/examples/ex_wdog.rs}

\lstinputlisting[style=ascii,caption={What \code{ex\_wdog} printed when the build ran it: the refresh, the warning, the reset, the two failures and the lock, and then the watchdog's Verilog.},label={lst:x-wdog-out}]{out_wdog.txt}
